\documentclass[a4paper,11pt]{article}
\usepackage[T1]{polski}
\usepackage[utf8]{inputenc}
\usepackage{xcolor}
\usepackage{titlesec}
\usepackage{graphicx}
\usepackage[inkscapeformat=pdf]{svg}
\usepackage{listings}
\usepackage{float}

\lstdefinestyle{custom}{
    backgroundcolor=\color{white},
    inputencoding=utf8,
    extendedchars=\true,
    breaklines=true,
    basicstyle=\ttfamily\footnotesize,
    tabsize=1
}
\lstset{style=custom}
\titlespacing*{\section}{0pt}{6pt}{3pt}
\titlespacing*{\subsection}{0pt}{4pt}{2pt}
\hoffset=-3.0cm
\textwidth=18cm
\evensidemargin=0pt
\voffset=-3cm
\textheight=27cm
\setlength{\parindent}{0pt}
\setlength{\parskip}{\medskipamount}
\raggedbottom
\newcommand\TY{\raise.17ex\hbox{$\scriptstyle\mathtt{\sim}$}}

\author{Matylda Drapich}
\title{Laboratorium sieci komputerowych\\ćw.6\\Serwer w chmurze, ZeroTier, kontener i SMB}
\frenchspacing

\begin{document}
\thispagestyle{empty}
\maketitle

\vspace{-0.5em}
Celem ćwiczenia była konfiguracja zapory sieciowej (Windows/Unix),
podłączenie serwera VPS w chmurze Azure do sieci ZeroTier ADM,
uruchomienie kontenera OCI (Docker) z usługą WWW oraz montaż
udziału SMB z laptopa. Pracę wykonano na stacji \texttt{c2};
rekwizytem był dostęp do serwera w chmurze przez sieć ADM
(\texttt{ssh \$srv}).

\section{Konfiguracja zapory internetowej}

\subsection{Unix -- Azure NSG (\texttt{LabSK-NSG})}

Skrypt \texttt{gen-vps} tworzy grupę bezpieczeństwa sieci w Azure
(\texttt{az/nsg -a}). Reguły NSG przepuszczają m.in.: SSH (22/TCP),
ICMP, ZeroTier (9993/UDP), HTTP (80/TCP) i HTTPS (443/TCP).

\textit{TODO: wkleić output \texttt{gen-vps} z fragmentem tworzenia NSG
(linie \texttt{Dodanie reguły AllowSSHInbound} itd.).}

\subsection{Windows / laptop -- udostępnienie SMB}

Na laptopie (macOS) utworzono katalog \texttt{\TY/vps} i włączono
udostępnianie plików (File Sharing) w celu eksportu katalogu dla VPS.
Na systemie Windows analogiczna konfiguracja wykonywana jest skryptem
\texttt{labsk/bin/Edit-smb.ps1} dla katalogu \texttt{c:\textbackslash vps}.

\textit{TODO: zrzut ekranu lub opis włączenia File Sharing na Macu
/ wynik \texttt{Edit-smb.ps1} na Windows.}

\section{Podłączenie serwera w chmurze do sieci ZeroTier ADM}

\subsection{Generacja VPS -- \texttt{gen-vps}}

Skrypt \texttt{gen-vps} rejestruje zasoby Azure, tworzy VM Archlinux,
zapisuje DNS serwera w \texttt{\TY/.vps} (polecenie \texttt{ssh \$(cat \TY/.vps)}),
ustawia zmienną \texttt{\$vps} w \texttt{\TY/.zshrc.local}
(polecenie \texttt{ssh \$vps}) oraz wpis \texttt{Host vps} w
\texttt{\TY/.ssh/config}. Po starcie VM dołącza serwer do sieci ADM.

\begin{lstlisting}
#!/bin/sh
test -f ~/.ssh/id_ed25519.pub || ssh-keygen -t ed25519 -f ~/.ssh/id_ed25519 -N ""
/pub/Linux/zetis/bin/az/register; /pub/Linux/zetis/bin/az/nsg -a
/pub/Linux/zetis/bin/az/archlinux -k; /pub/Linux/zetis/bin/az/archlinux -s
echo $USER.polandcentral.cloudapp.azure.com >~/.vps
echo 'vps=$(cat ~/.vps); export vps' >~/.zshrc.local
# ... ~/.ssh/config ...
ssh $(cat ~/.vps) <<'E'
sudo pacman -S --needed --noconfirm zerotier-one
sudo zerotier-cli join 8d1c312afa9a3adc
E
\end{lstlisting}

\textit{TODO: wkleić logi: \texttt{gen-vps}, \texttt{cat \TY/.vps},
\texttt{ssh \$(cat \TY/.vps) hostname}, \texttt{ssh \$vps hostname}.}

\subsection{Konfiguracja VPS -- \texttt{config-vps}}

Skrypt \texttt{config-vps} instaluje na VPS m.in.\ \texttt{zsh},
\texttt{zerotier-one}, \texttt{gemini-cli} oraz tworzy narzędzie \texttt{ipb}.

\begin{lstlisting}
#!/bin/sh
sleep 15
ssh -T vps <<'E'
sudo pacman -Syy --noconfirm
sudo pacman -S --needed --noconfirm zsh zerotier-one gemini-cli
sudo systemctl enable --now zerotier-one.service
sudo usermod -s /usr/bin/zsh $USER
sudo zerotier-cli join 8d1c312afa9a3adc
E
\end{lstlisting}

\textit{TODO: wkleić \texttt{config-vps}, \texttt{ssh vps "ipb a"} z interfejsem
\texttt{10.80.108.x/24}, \texttt{listnetworks OK}.}

\subsection{Sieć ADM -- parametry i schemat}

Sieć ZeroTier ADM: NetID \texttt{8d1c312afa9a3adc},
zakres \texttt{10.80.108.0/24}. Urządzenia: VPS (\texttt{vps.adm}),
laptop (\texttt{lap.adm}), stacja uczelniana (\texttt{s3} / \texttt{c2}).

\begin{figure}[H]
  \centering
  \includesvg[width=0.58\textwidth]{ADM.svg}
  \caption{Schemat sieci ADM}
\end{figure}

\textit{TODO: ping \texttt{vps.adm} / \texttt{lap.adm} z MacBooka lub c2.}

\section{Kontener OCI (Docker) i usługa WWW}

\subsection{Skrypt \texttt{gen-kont}}

Skrypt \texttt{gen-kont} (na stacji \texttt{c2}) zdalnie instaluje Docker
na VPS, uruchamia kontener \texttt{nginx} z mapowaniem portu
\texttt{8080:80} i weryfikuje dostęp publicznie oraz przez ADM.

\begin{lstlisting}
#!/bin/sh
. ~/.zshrc.local
ssh $vps <<'E'
sudo pacman -S --needed --noconfirm docker
sudo systemctl enable --now docker
sudo docker run -d --name web -p 8080:80 nginx
E
curl -sI http://$vps:8080 | head -1
curl -sI http://vps.adm:8080 | head -1
\end{lstlisting}

\textit{TODO: wkleić wynik \texttt{gen-kont} --- dwa wiersze
\texttt{HTTP/1.1 200 OK} (adres publiczny i \texttt{vps.adm}).}

\subsection{Konfiguracja sieci kontenera}

Kontener nasłuchuje na porcie \texttt{8080} hosta VPS; ruch jest
kierowany przez reguły NSG (TCP 8080 --- \textit{TODO: jeśli wymagane,
dodać regułę lub użyć portu 80}) oraz przez interfejs ZeroTier
(\texttt{vps.adm:8080}).

\textit{TODO: \texttt{ssh vps "sudo docker ps"} oraz
\texttt{curl -I http://vps.adm:8080}.}

\section{Usługa SMB -- \texttt{mount-smb}}

Na laptopie udostępniono katalog \texttt{\TY/vps}. Na serwerze VPS
skrypt \texttt{mount-smb} montuje udział \texttt{//lap.adm/vps}
w katalogu \texttt{\TY/vps}.

\begin{lstlisting}
#!/bin/sh
d=~/vps; mkdir -p $d
sudo pacman -S --needed --noconfirm cifs-utils
sudo mount -t cifs //lap.adm/vps $d -o user=$SMB_USER,pass=$SMB_PASS,uid=$(id -u),vers=3.0
ls $d
\end{lstlisting}

\textit{TODO: wkleić \texttt{export SMB\_USER=... SMB\_PASS=...},
\texttt{mount-smb}, wynik \texttt{ls \TY/vps} z plikiem testowym z laptopa.}

\section{Podsumowanie}

\begin{itemize}
\item \textbf{Zapora:} NSG w Azure (Unix/chmura); udostępnienie SMB na laptopie.
\item \textbf{ZeroTier ADM:} VPS dołączony do sieci \texttt{8d1c312afa9a3adc}.
\item \textbf{Kontener:} Docker + nginx, dostęp przez \texttt{\$vps} i \texttt{vps.adm}.
\item \textbf{SMB:} montaż \texttt{\TY/vps} na VPS z laptopa.
\end{itemize}

\end{document}
